\chapter{Detección Temprana y Artefactos en Vivo}

La respuesta a incidentes informáticos comienza invariablemente con una anomalía: una discrepancia operativa, un proceso con consumo atípico o un archivo fuera de su ubicación canónica. A mediados de abril de 2004, la administración de cómputo del instituto descubrió que el servidor principal departamental (\texttt{victima1}) albergaba ejecutables no autorizados configurados con privilegios de superusuario (\texttt{root}). 

Este capítulo documenta las primeras acciones de respuesta en vivo (\textit{live response}), la metodología de búsqueda exhaustiva en el almacenamiento masivo y el análisis forense de los artefactos inmediatos que revelaron el compromiso inicial del sistema.

\begin{notaforense}[Uso Tradicional de /dev como Mecanismo de Ocultamiento]
Históricamente, el directorio \texttt{/dev} constituyó el refugio predilecto de los intrusos en entornos Unix y GNU/Linux. Debido a que este directorio aloja cientos de nodos de dispositivos de bloques y caracteres legítimos (como \texttt{/dev/null}, \texttt{/dev/tty} o \texttt{/dev/hda}), los atacantes solían crear subdirectorios camuflados con nombres que simulaban interfaces virtuales (por ejemplo, \texttt{/dev/ptyp} o \texttt{/dev/zqft4}) para depositar código malicioso, herramientas de compilación y \textit{sniffers}, asumiendo que los administradores rara vez inspeccionaban minuciosamente esta jerarquía.
\end{notaforense}

\section{Detección Inicial: Archivos con Permisos SUID y SGID No Autorizados}

El descubrimiento que detonó la investigación ocurrió durante una revisión de mantenimiento rutinario. Al detectar inconsistencias en los permisos de ciertos programas, el equipo pericial procedió a ejecutar una auditoría integral sobre el árbol de directorios para identificar la totalidad de los archivos con los bits \textbf{SUID} (\textit{Set User ID}) y \textbf{SGID} (\textit{Set Group ID}) activados.

Como principio de arquitectura en sistemas Unix, un binario con el bit SUID activo se ejecuta con los privilegios del propietario del archivo y no con los del usuario que lo invoca. Si un archivo perteneciente a \texttt{root} posee el bit SUID habilitado (\texttt{-rwsr-xr-x}), cualquier usuario local sin privilegios que lo invoque adquirirá transitoriamente las facultades del superusuario.

El barrido inicial reveló la presencia de dos ejecutables altamente sospechosos alojados en el directorio personal del usuario \texttt{duende} (UID 500)\footnote{Por razones de confidencialidad institucional, en este libro se emplean seudónimos técnicos estandarizados para identificar a las cuentas de usuario, preservando con exactitud los identificadores numéricos (UID/GID) y las rutas internas.}, tal como se documenta en el Listado~\ref{lst:suid_duende}.

\begin{lstlisting}[style=terminal,caption={Ejecutables con bit SUID activo localizados en el directorio del usuario duende.},label={lst:suid_duende}]
# find / -type f \( -perm -04000 -o -perm -02000 \) -ls
702242   92 -rwsr-sr-x   1 root     root        88531 Feb 11  1999 /home/staff/duende/prog/bin/registro
620791   20 -rwsrwsr-x   1 root     root        18256 Jan  9  2000 /home/staff/duende/prog/gtk/registros
\end{lstlisting}

Un análisis preliminar de estos dos binarios confirmó que se trataba de programas compilados cuyos propósitos funcionales no justificaban bajo ninguna circunstancia privilegios de superusuario. Más crítico aún: aunque residían en la cuenta de un usuario ordinario, su propietario formal había sido alterado a \texttt{root:root} con permisos de lectura, escritura y ejecución universal (\texttt{04755} y \texttt{02775}). 

Aunque la administración procedió de inmediato a revocar los permisos SUID para mitigar el riesgo de elevación de privilegios inmediata, la gravedad del hallazgo exigió la apertura formal de una investigación pericial profunda.

\begin{retoforense}[Reto Pericial 1: Auditoría de Permisos SUID y Triage en Vivo]
Suponga que durante una respuesta a incidentes en un servidor GNU/Linux localiza un ejecutable en \texttt{/tmp/run} con permisos \texttt{-rwsr-xr-x} perteneciente a \texttt{root}:
\begin{enumerate}
    \item ¿Qué comandos de inspección binaria utilizaría para examinar el archivo de forma pasiva sin arriesgar la ejecución de código malicioso?
    \item ¿Cómo determinaría mediante el comando \texttt{stat} el número de inodo y si el archivo mantiene enlaces duros (\textit{hard links}) en otras ramas del disco?
    \item Si revoca preventivamente el bit especial con \texttt{chmod u-s /tmp/run}, ¿cuál de las tres marcas temporales del inodo (\textit{mtime}, \textit{atime} o \textit{ctime}) se actualizará en ese instante y cuál permanecerá intacta?
\end{enumerate}
\end{retoforense}


\section{Presencia Anómala en Directorios del Sistema}

Tras confirmar la anomalía en el espacio de usuario, la siguiente fase pericial consistió en evaluar la integridad de los directorios del sistema (\texttt{/usr/bin}, \texttt{/bin}, \texttt{/sbin}, \texttt{/dev}). Para ello se ejecutó una búsqueda diferencial de todos los objetos en el disco duro cuyo propietario fuera un identificador distinto de \texttt{root} (UID 0), descartando los directorios personales legítimos montados en \texttt{/home}.

El resultado de esta inspección arrojó una violación flagrante del principio de segregación: el usuario \texttt{duende} (UID 500) figuraba como propietario directo de utilidades críticas del sistema operativo y de scripts encubiertos bajo \texttt{/dev}, como se ilustra en el Listado~\ref{lst:archivos_sistema_duende}.

\begin{lstlisting}[style=terminal,caption={Archivos en directorios del sistema pertenecientes al usuario duende (UID 500).},label={lst:archivos_sistema_duende}]
501289   24 -rwxr-xr-x   1 500      users       22837 Apr 11 16:18 ./usr/bin/top
501292   24 -rwxr-xr-x   1 500      users       22837 Apr 11 16:19 ./usr/bin/pstree
388613   24 -rwxr-xr-x   1 500      users       20705 Apr 11 16:24 ./usr/bin/getty
130454    4 -rwxr-xr-x   1 500      users          36 Nov 14  2002 ./dev/.libn
130456    4 -rwxr-xr-x   1 500      users          66 Mar  8 06:12 ./dev/.libpt
130461    4 -rwxr-xr-x   1 500      users          69 Mar  8 06:13 ./dev/.libtr
130458    4 -rwxr-xr-x   1 500      users          65 Mar  8 06:12 ./dev/.libtp
533480    4 drwxr-xr-x   2 500      users        4096 May  9 19:05 ./dev/zqft4
\end{lstlisting}

\subsection{Evaluación Forense de los Binarios Alterados}

El hallazgo de los archivos listados reviste una severidad técnica extraordinaria:
\begin{enumerate}
    \item \textbf{Herramientas de monitorización de procesos (\texttt{top} y \texttt{pstree}):}
    Los comandos \texttt{/usr/bin/top} y \texttt{/usr/bin/pstree} son componentes fundamentales para la supervisión operativa del servidor. Al pertenecer al usuario 500 y exhibir fechas de modificación recientes (11 de abril a las 16:18~h), surgió la certeza de que habían sido suplantados por versiones troyanizadas diseñadas para omitir u ocultar los procesos maliciosos del intruso ante la vista del administrador.
    \item \textbf{Gestión de terminales (\texttt{getty}):}
    El ejecutable \texttt{getty} gestiona los terminales de acceso interactivo y solicitudes de autenticación en consola física o serie. Su alteración apuntaba directamente a la captura encubierta de contraseñas de inicio de sesión.
    \item \textbf{Estructura clandestina en \texttt{/dev}:}
    Los cuatro archivos con prefijo oculto (\texttt{.libn}, \texttt{.libpt}, \texttt{.libtr}, \texttt{.libtp}) y el subdirectorio anómalo \texttt{/dev/zqft4} confirmaron que el intruso había establecido un espacio de trabajo persistente dentro del directorio de dispositivos.
\end{enumerate}

Para descartar la intervención de otros usuarios de la institución, se repitió la consulta para cada cuenta registrada en el sistema. Los resultados confirmaron que ningún otro usuario presentaba archivos fuera de su directorio personal, focalizando la línea de investigación sobre la cuenta \texttt{duende}.

\section{Auditoría de Bitácoras de Intérprete y Técnicas Antiforense}

Para reconstruir las acciones ejecutadas por el usuario \texttt{duende}, el procedimiento estándar exigía examinar los archivos de historial de comandos generados por las \textit{shells} interactivas: \texttt{.bash\_history} y \texttt{.history}.

Al inspeccionar el directorio personal de la cuenta sospechosa se registraron las siguientes condiciones:
\begin{itemize}
    \item El archivo \texttt{.bash\_history} se encontraba completamente vacío (tamaño 0 bytes).
    \item El archivo \texttt{.history} (correspondiente al intérprete \texttt{tcsh}) contenía únicamente comandos antiguos e inconsecuentes correspondientes a meses anteriores.
\end{itemize}

\begin{advertencia}[Indicio de Técnicas Antiforense]
En un entorno de producción, la ausencia total de registros en el historial de comandos de un usuario con actividad recurrente constituye un indicador clásico de técnicas antiforense (\textit{anti-forensics}). El vaciado deliberado de bitácoras busca entorpecer la reconstrucción cronológica y ocultar los parámetros utilizados en la invocación de herramientas maliciosas.
\end{advertencia}

\section{Descubrimiento de Trazas Ocultas: El Registro de Midnight Commander}

Ante la ausencia de registros en los historiales estándar, el equipo pericial recurrió a fuentes de telemetría secundaria comúnmente ignoradas por los atacantes. Una de las aplicaciones interactivas de gestión de archivos más extendidas en la época era \textit{Midnight Commander} (\texttt{mc}), un clon de Norton Commander para terminales Unix.

A diferencia de las \textit{shells} convencionales, \texttt{mc} almacena de manera silenciosa el historial de directorios visitados y archivos abiertos en el archivo interno \texttt{.mc/history}, ubicado dentro del perfil de cada usuario. 

La extracción y análisis del archivo \texttt{/home/staff/duende/.mc/history} permitió recuperar la secuencia exacta de exploración realizada por el intruso durante el incidente, como se documenta en el Listado~\ref{lst:mc_history}.

\begin{lstlisting}[style=terminal,caption={Extracto de navegación recuperado del historial de Midnight Commander (.mc/history).},label={lst:mc_history}]
[Dir Hist]
0=/home/staff/duende/prog
1=/home/staff/duende/prog/c
2=/home/staff/duende/prog/c/uno
3=/home/staff/duende/prog/c/dos
4=/home/staff/duende
5=/home/staff/duende/bin
6=/home/staff/duende
7=/home/staff/duende/prog
8=/home/staff/duende/prog/bin
9=/home/staff/duende/prog
10=/home/staff/duende
28=/root/archivos
46=cd /dev
47=/dev
48=/dev/zqft4
49=/dev
50=/home/staff/duende
51=cd /dev
52=/dev
53=/home/staff/duende
54=/dev
55=/dev/zqft4
56=/dev
57=/home/staff/duende
\end{lstlisting}

\subsection{Interpretación Pericial de la Navegación}

El análisis cronológico de los registros de \texttt{mc} reveló conclusiones probatorias determinantes:
\begin{enumerate}
    \item \textbf{Exploración recurrente de \texttt{/dev/zqft4}:} El usuario navegó repetidamente entre su directorio personal y el directorio clandestino \texttt{/dev/zqft4}, desvirtuando cualquier argumento de desconocimiento involuntario.
    \item \textbf{Acceso a directorios restringidos de administración:} La entrada número 28 demuestra un acceso directo a \texttt{/root/archivos}, una ruta accesible únicamente mediante privilegios de superusuario o brechas en la configuración de permisos.
    \item \textbf{Manipulación de directorios de correo y cuentas ajenas:} En las secciones subsiguientes del registro de \texttt{mc} se constató que la cuenta intentó visualizar archivos gráficos pertenecientes a otros usuarios (\texttt{brenton}) y exploró la cola de correo en \texttt{/var/mail}.
\end{enumerate}

\section{Conclusiones de la Fase de Detección Temprana}

Las evidencias materiales recabadas en esta fase inicial demostraron con certeza pericial que:
\begin{itemize}
    \item El servidor \texttt{victima1} se encontraba efectivamente comprometido a nivel de superusuario.
    \item Se habían sembrado binarios con privilegios SUID y alterado utilidades de administración esenciales (\texttt{top}, \texttt{pstree}, \texttt{getty}).
    \item Existía un espacio de almacenamiento malicioso encubierto bajo \texttt{/dev}.
    \item La cuenta \texttt{duende} (UID 500) mantenía un vínculo material directo con los ejecutables y directorios ilícitos.
\end{itemize}

Frente a este escenario, la administración tomó la determinación táctica de no interrumpir de forma inmediata la operación del equipo, sino instaurar un dispositivo de \textbf{monitoreo encubierto y pasivo} para identificar los vectores de comunicación externa y esclarecer si se trataba de un ataque automatizado foráneo o de una amenaza interna deliberada.
